\scanpage{23}
- 218 - \hfill n\textsuperscript{o} 391

(VII) La demi-somme des racines positives est la somme des poids fondamentaux, (§ 6, n\textsuperscript{o} 9, prop. 8), et vaut donc :
\[
\begin{aligned}
\rho&=\varepsilon_2+2\varepsilon_3+3\varepsilon_4+4\varepsilon_5+4(\varepsilon_8-\varepsilon_7-\varepsilon_6)=\\
&=8\alpha_1+11\alpha_2+15\alpha_3+21\alpha_4+15\alpha_5+8\alpha_6.
\end{aligned}
\]

(VIII) De même qu'au n\textsuperscript{o} 11, le réseau des vocables $\Gamma_{\unsure{o|0}}(R'')$ est l'associé \uline{dans} $V''$ du réseau $\Gamma_2^*$ de $V''{}^*$ \uline{projection orthogonale} dans $V''{}^*$ de $\Gamma_{\unsure{o|0}}(R^*)$, (lorsque $V''{}^*$ est identifié au sous-espace de $V^*$ orthogonal à $\varepsilon_7^*+\varepsilon_8^*$ et $\varepsilon_7^*-\varepsilon_6^*$). Le groupe $\Gamma(R'')/\Gamma_{\unsure{o|0}}(R'')$ est donc \uline{isomorphe au groupe} $\Gamma_2^*/\Gamma_{\unsure{o|0}}(R''{}^*)$, qu'il reste à calculer. Or, les projections orthogonales de $\varepsilon_8^*$, $-\varepsilon_7^*$ et $-\varepsilon_6^*$ sur $V''{}^*$ sont toutes trois égales $\frac13(\varepsilon_8^*-\varepsilon_7^*-\varepsilon_6^*)$ et les autres $\varepsilon_i^*$ sont dans $V''{}^*$ ; on en déduit aisément que $\Gamma_2^*$ est engendré par $\Gamma_{\unsure{o|0}}(R''{}^*)$ et par l'unique vecteur $\mu^*=\frac13(\varepsilon_8^*-\varepsilon_7^*-\varepsilon_6^*)$ ; mais comme $3\mu^*\in\Gamma_{\unsure{o|0}}(R''{}^*)$, on voit que $\Gamma_2^*/\Gamma_{\unsure{o|0}}(R''{}^*)$ est isomorphe à $\mathrm Z/3\mathrm Z$ et par suite $\Gamma(R'')/\Gamma_{\unsure{o|0}}(R'')$ est \uline{isomorphe à} $\mathrm Z/3\mathrm Z$, et l'indice de connexion $f=3$.

(IX) et (X). Le calcul de la plus grande racine et de l'indice de connexion permet, grâce au § 6, n\textsuperscript{o} 10, prop. 23, de calculer l'ordre du groupe de Weyl $W(R'')$, qui est égal à :
\[
6!1.2.2.3.2.1.3=2^7.3^4.5=51.840.
\]
La suite des exposants a 6 termes compris entre 1 et 11, et comprend les termes 1, 5, 7, 11 qui sont étrangers à 12, (§ 8, n\textsuperscript{o} 5, prop. 6). Les deux derniers termes de la suite $m$, $m'$ sont des entiers tels que :
\[
m+m'=12
\]
\[
(m+1)(m'+1)(1+1)(5+1)(7+1)(12+1)=2^7.3^4.5
\]
\mspage{22}
